\subsection{Multi-Modal Learning}
% Multi-modal learning has made significant progress on image-text tasks recent year, especially on large scale pre-training tasks, with amazing results.Their paradigm is mainly to use two encoders to respectively encode image and text information, and use contrastive learning for training. CLIP~\citep{radford2021learning} is a milestone work. It uses 400 million image-text data crawled on the Internet for training, and has achieved remarkable results in image-text retrieval task, zero-shot image recognition tasks and pre-training task. ALIGN~\citep{jia2021scaling} scales up the image-text data to 1 billion on the basis of CLIP~\citep{radford2021learning}. During training, the text encoder generates the “label” weights and uses a larger batchsize to achieve better results than CLIP~\citep{radford2021learning}.WenLan~\citep{huo2021wenlan} and EfficientCLIP~\citep{wang2021efficientclip} train Chinese pairs data. EfficientCLIP~\citep{wang2021efficientclip} uses Ensemble Confident Learning to make good use of single-modal text data during multi-modal training. With fewer training resources, it achieves the state-of-the-art performance on Chinese cross-modal retrieval tasks. CLIP2Video~\citep{fang2021clip2video} is based on the semantic space constructed by the CLIP~\citep{radford2021learning} model to capture motions at fine temporal video frames, re-align the tokens of video clips and phrases and enhance the multi-modal correlation to transfer the image-language pre-training model to video-text retrieval, achieve state-of-the-art performance on major text-to-video and video-to-text retrieval benchmarks.

% Although these works have achieved significant results on the basis of contrastive learning, because only the contrastive learning between image and text is carried out, the self-supervision of image and text  is ignored, so that when training large scale data, there is low data utilization. Therefore, we think more comprehensively about vision-language contrastive learning, and form a group contrastive learning  paradigm  composed of Image-Image contrastive learning, Text-Text contrastive learning, Image-Text contrastive learning, and Nearest-Neighbor-Text contrastive learning. 